\documentclass[10pt,a4paper]{amsart}
\usepackage[margin=19mm]{geometry}
\usepackage{amsmath,amssymb,mathtools}
\usepackage[scr=rsfs,cal=euler]{mathalfa}
\usepackage{xfrac}
\usepackage[unicode,hidelinks,bookmarks=false]{hyperref}
\newcommand{\ie}{i.e.\ }
\newcommand{\cf}{cf.\ }
\newcommand{\resp}{respectively\ }
\newcommand{\Subsection}{\subsection}
\providecommand{\nobreakdash}{\nobreak\hbox{-}}
\setlength{\emergencystretch}{2em}
\multlinegap0pt
\usepackage{bm}
\usepackage{bbm}
\usepackage{dsfont}
\usepackage{xspace}
\usepackage{cancel}
\usepackage{mathtools}
\usepackage{amsthm}
\makeatletter
\@ifundefined{lemma}{\newtheorem{lemma}{Lemma}}{}
\makeatother
\title[Stabilization of arbitrary structures in a three-dimensional]{Stabilization of arbitrary structures in a three-dimensional doubly degenerate nutrient taxis system}
\author{De-Ji-Xiang-Mao, Ai Huang, Yifu Wang}
\date{Source: arXiv:2601.12218v1 (CC BY 4.0)}
\begin{document}
\maketitle

\noindent\emph{Source and licence.} Stabilization of arbitrary structures in a three-dimensional doubly degenerate nutrient taxis system. Version arXiv:2601.12218v1 (CC BY 4.0), distributed under the Creative Commons Attribution 4.0 International licence, as recorded in the arXiv licence metadata for this exact version and shown on the abstract page https://arxiv.org/abs/2601.12218v1. Full rights notice: licenses/arxiv-2601.12218-CC-BY-4.0.txt.\par\smallskip

\noindent\textit{Editorial scope.} Counted unit: the Lemma whose source label is \texttt{lemma3.8} in the article (source0.tex lines 1183--1195, source label \texttt{lemma3.8}), delivered together with the article's own complete original proof of it (source0.tex lines 1196--1276, one \texttt{proof} environment of 81 lines). No mathematics is written, repaired or re-derived: statement and proof are the article's own text line for line. Required context retained uncounted: 1.6 (lines 217--219); 2.5 (lines 380--390); 2.10 (lines 422--424); 2.11 (lines 436--438); 2.18 (lines 522--524); lemma2.7 (lines 652--666); 2.29 (lines 655--659); 3.15 (lines 1051--1053); lemma3.7 (lines 1118--1129). Every definition, display and label of those lines is retained unchanged. Omitted from this package and not redistributed: 1--216, 220--379, 391--421, 425--435, 439--521, 525--651, 660--1050, 1054--1117, 1130--1182, 1277--1741. Nothing inside a retained excerpt is omitted or altered: every retained body is delivered exactly as the article's lines are written, so no omission receipt of any kind accompanies this package. Citations: the retained excerpts carry no citation command at all, so no bibliography entry of the article is transcribed and no reference is redistributed. Reference adaptation: none was needed and none was made; every reference inside the delivered unit resolves inside it, so no unresolved reference or citation is produced. Printed slips kept literal: no printed word, symbol, hypothesis or number of the counted statement or of its proof was altered. Layout and class: the article's own class is \texttt{article} with packages including amssymb, amsmath, amsthm, amsfonts, graphicx, epsfig, tikz; the wrapper is the lane's pinned \texttt{amsart} editorial wrapper at 10pt on A4, which restates the theorem-like environments the retained text needs and adds the article's own macro definitions that the retained text needs (none), with extra wrapper packages (bm, bbm, dsfont, xspace, cancel, mathtools). The delivered numbering is the unit's own. Assets: no retained span contains a figure environment, a graphics-inclusion command, a PDF-page inclusion, a table or a TikZ picture; no image file is copied into this package, the package declares no asset dependency and no image file is redistributed with it. Licence: version arXiv:2601.12218v1 of this article is distributed under the Creative Commons Attribution 4.0 International licence, as recorded in the arXiv licence metadata for this exact version and shown on the abstract page https://arxiv.org/abs/2601.12218v1. A full rights notice is delivered at licenses/arxiv-2601.12218-CC-BY-4.0.txt. Characters: every retained excerpt is pure ASCII, so the delivered engine, XeLaTeX with its default Latin Modern text font, reports no missing character.
\begin{align}\label{1.6}
    \|u_0\|_{L^{\infty}(\Omega)}+\|v_0\|_{L^{\infty}(\Omega)}+\|\nabla \ln v_0\|_{L^{\infty}(\Omega)}\leq K\quad \text{and}\quad \int_{\Omega}\ln{u_0}\geq -K
\end{align}

\begin{equation}\label{2.5}
\left\{
\begin{aligned}
&u_{\varepsilon t}=\nabla \cdot (u_{\varepsilon}v_{\varepsilon}\nabla u_{\varepsilon})-\chi \nabla \cdot
(u^{\alpha}_{\varepsilon}v_{\varepsilon}\nabla v_{\varepsilon})+\ell u_{\varepsilon}v_{\varepsilon},&x\in \Omega,\, t>0,\\
& v_{\varepsilon t}=\Delta v_{\varepsilon}-u_{\varepsilon}v_{\varepsilon},&x\in \Omega,\, t>0,\\
& \frac{\partial u_{\varepsilon }}{\partial \nu}=\frac{\partial v_{\varepsilon }}{\partial \nu}=0,&x\in \partial\Omega,\, t>0,\\
& u_{\varepsilon }(x,0)=u_0(x)+\varepsilon,\;v_{\varepsilon }(x,0)=v_0(x),&x\in \Omega,\,
\end{aligned}
\right.
\end{equation}

\begin{align}\label{2.10}
    \int_{0}^{T_{max,\varepsilon}}\int_{\Omega}u_{\varepsilon }v_{\varepsilon }\leq \int_{\Omega}v_{0}.
\end{align}

\begin{align}\label{2.11}
    \int_0^{T_{max,\varepsilon}}\int_{\Omega}v_{\varepsilon}|\nabla v_{\varepsilon}|^2\leq C.
\end{align}

\begin{align}\label{2.18}
\int_0^{T_{\max, \varepsilon}}\int_{\Omega}v_{\varepsilon}\leq C
\end{align}

\begin{lemma}\label{lemma2.7}
Let $\Omega \subset \mathbb R^3$ be a bounded domain with smooth boundary and $q\geq 2$. Then for all $t\in (0,T_{max,\epsilon})$, there exist
$\Gamma(q)>0$ and $\gamma(q)>0$ such that
\begin{align}\label{2.29}
    \frac{d}{dt}\int_{\Omega}\frac{|\nabla v_{\epsilon}|^q}{v_{\epsilon}^{q-1}}+\Gamma(q)\int_{\Omega}\frac{|\nabla
    v_{\epsilon}|^{q-2}}{v_{\epsilon}^{q-3}}|D^2\ln{v_{\epsilon}}|^2+\Gamma(q)\int_{\Omega}u_{\epsilon}\frac{|\nabla
    v_{\epsilon}|^q}{v_{\epsilon}^{q-1}}\leq \frac{1}{\Gamma(q)}\int_{\Omega}v_{\epsilon}(u_{\epsilon}^{\frac{q+2}{2}}+1)
\end{align}
and
\begin{align}\label{2.30}
    \frac{d}{dt}\int_{\Omega}\frac{|\nabla v_{\varepsilon}|^q}{v_{\varepsilon}^{q-1}}+\gamma(q)\int_{\Omega}\frac{|\nabla
    v_{\varepsilon}|^{q-2}}{v_{\varepsilon}^{q-3}}|D^2\ln{v_{\varepsilon}}|^2+\gamma(q)\int_{\Omega}u_{\varepsilon}\frac{|\nabla
    v_{\varepsilon}|^q}{v_{\varepsilon}^{q-1}}\leq \frac{1}{\gamma(q)}\int_{\Omega}v_{\varepsilon}(|\nabla u_{\varepsilon}|^{\frac{q+2}{3}}+1).
\end{align}
\end{lemma}

\begin{align}\label{3.15}
    \int_{\Omega}u_{\varepsilon}^{p_0}(\cdot,t)\leq C\qquad for\;all\; t\in (0,T_{max,\varepsilon}).
\end{align}

\begin{lemma}\label{lemma3.7}
%Let $\alpha\in(\frac{4}{3},\frac{19}{12})$, $p>1$ and  $K>0$ with the
%property that \eqref{1.5}  holds, and  assume that $T_{max,\varepsilon}<\infty$.
For $L>0,p_0>\frac{3}{2}$, $p>1$, $q>2+\frac{3p}{p_0}$, $\beta\in [1,\frac{2p_0}{3}+p+1)$
and any $\eta\in(0,1)$, there exists $C=C(\eta, L,p,\beta
 )>0$ such that
whenever  $\int_{\Omega}\varphi^{p_0}\leq L$,
  \begin{align}\label{3.25}
    \int_{\Omega}\varphi^{\beta} \psi\leq \eta \int_{\Omega}\varphi^{p-1}\psi|\nabla
    \varphi|^2+\eta \int_{\Omega}\varphi\frac{|\nabla  \psi|^q}{{\psi}^{q-1}}+C\int_{\Omega}\varphi \psi
\end{align} holds for   all $\varphi\in C^1(\overline \Omega)$ and $\psi \in C^1(\overline \Omega)$ fulfilling $\varphi> 0$ and $\psi>0$ in $\overline \Omega$.
\end{lemma}

\begin{lemma}\label{lemma3.8}
Let $\Omega \subset \mathbb R^3$ and $K>0$ with the property that \eqref{1.6} holds. Supposed that
$\alpha\in\left(\frac{3}{2},\frac{19}{12}\right)$  and $p>1$, there exists  $C=C( K, p)>0$ such that for all
 $ t\in (0,T_{max,\varepsilon})$,
\begin{align}\label{3.30}
   \int_{\Omega}u_{\varepsilon}^{p}\leq C\quad \hbox{and}\quad
   \int_0^{T_{max,\varepsilon}}\int_{\Omega}u_{\varepsilon}^{p-1}v_{\varepsilon}|\nabla u_{\varepsilon}|^2\leq C
\end{align}
as well as
\begin{align}\label{3.30-1}
   \int_0^{T_{max,\varepsilon}}\int_{\Omega}u_{\varepsilon}^{p+1}v_{\varepsilon}\leq C.
\end{align}
\end{lemma}

\begin{proof}
From \eqref{2.29} in Lemma \ref{lemma2.7}, it follows that
    \begin{align}\label{3.31}
    \frac{d}{dt}\int_{\Omega}\frac{|\nabla v_{\varepsilon}|^q}{v_{\varepsilon}^{q-1}}+\Gamma(q)\int_{\Omega}\frac{|\nabla
    v_{\varepsilon}|^{q-2}}{v_{\varepsilon}^{q-3}}|D^2\ln{v_{\varepsilon}}|^2+\Gamma(q)\int_{\Omega}u_{\varepsilon}\frac{|\nabla
    v_{\varepsilon}|^q}{v_{\varepsilon}^{q-1}}\leq \frac{1}{\Gamma(q)}\int_{\Omega}v_{\varepsilon}(u_{\varepsilon}^{\frac{q+2}{2}}+1)
\end{align}
for all $t\in (0,T_{max,\varepsilon})$.

  Multiplying the first equation in \eqref{2.5} by $u_{\varepsilon}^{p-1}$, applying Young's inequality and \eqref{1.6}, there exists
  $c_1>0$ such that
    \begin{align}\label{3.32}
      &\frac{1}p \frac{d}{dt}\int_{\Omega}u_{\varepsilon}^{p}+\frac{p-1}{2}\int_{\Omega}u_{\varepsilon}^{p-1}v_{\varepsilon}|\nabla
      u_{\varepsilon}|^2\nonumber
      \\
      &\leq \frac{p-1}{2}\int_{\Omega}u_{\varepsilon}^{p+2\alpha-3}v_{\varepsilon}|\nabla
      v_{\epsilon}|^2+\ell\int_{\Omega}u_{\varepsilon}^pv_{\varepsilon}
      \\
      &\leq (p-1)\int_{\Omega}u_{\varepsilon}^{p+1}v_{\varepsilon}|\nabla v_{\varepsilon}|^2+(p-1)\int_{\Omega}v_{\varepsilon}|\nabla
      v_{\epsilon}|^2+\ell\int_{\Omega}u_{\varepsilon}^{p+1}v_{\varepsilon}+\ell\int_{\Omega}u_{\varepsilon}v_{\varepsilon}\nonumber
      \\
      &\leq  \frac{\Gamma(q)} 2
      \int_{\Omega}u_{\varepsilon}\frac{|\nabla v_{\varepsilon}|^q}{v_{\varepsilon}^{q-1}}+(p-1)\int_{\Omega}v_{\varepsilon}|\nabla
      v_{\varepsilon}|^2+\ell\int_{\Omega}u_{\varepsilon}^{p+1}v_{\varepsilon}+c_1\int_{\Omega}u_{\varepsilon}^{\frac{qp+q-2}{q-2}}v_{\varepsilon}^{\frac{3q-2}{q-2}}+\ell\int_{\Omega}u_{\varepsilon}v_{\varepsilon}\nonumber
      \\
      &\leq    \frac{\Gamma(q)} 2\int_{\Omega}u_{\varepsilon}\frac{|\nabla v_{\varepsilon}|^q}{v_{\varepsilon}^{q-1}}+(p-1)\int_{\Omega}v_{\varepsilon}|\nabla
      v_{\varepsilon}|^2+\ell\int_{\Omega}u_{\varepsilon}^{p+1}v_{\varepsilon}+
      c_1K^{\frac{2q}{q-2}}\int_{\Omega}u_{\varepsilon}^{\frac{qp+q-2}{q-2}}v_{\varepsilon}+
      \ell\int_{\Omega}u_{\varepsilon}v_{\varepsilon}\nonumber
    \end{align}
    for all $t\in (0,T_{max,\varepsilon})$.

     Noting that  if $p_0>\frac{3}{2}$, then for any  fixed $p>1$, one can find $q>2$ satisfying
\begin{align}\label{3.34}
    q\in \left( 2+\frac{3p}{p_0},\frac{4p_0}{3}+2p \right)
\end{align}
which ensures that
\begin{align}\label{3.35}
    \frac{qp+q-2}{q-2}<\frac{2p_0}{3}+p+1
\end{align}
and
\begin{align}\label{3.36}
  \frac{q+2}{2}<\frac{2p_0}{3}+p+1.
\end{align}
Therefore thanks to \eqref{3.15}, \eqref{3.35} and \eqref{3.36}, the application of Lemma \ref{lemma3.7} to $\eta:=\min\{\frac{\Gamma(q)} 4, \frac{p-1}4\}$  yields
\begin{align}\label{3.37}
   &\ell\int_{\Omega}u_{\varepsilon}^{p+1}v_{\varepsilon}+c_1K^{\frac{2q}{q-2}}\int_{\Omega}u_{\varepsilon}^{\frac{qp+q-2}{q-2}}v_{\varepsilon}+\frac{1}{\Gamma(q)}\int_{\Omega}u_{\varepsilon}^{\frac{q+2}{2}}v_{\varepsilon}\nonumber
   \\
   &\leq\eta\int_{\Omega}u_{\varepsilon}^{p-1}v_{\varepsilon}|\nabla u_{\varepsilon}|^2+\eta\int_{\Omega}u_{\varepsilon}\frac{|\nabla
   v_{\varepsilon}|^q}{v_{\varepsilon}^{q-1}}+C(\eta,K,p)\int_{\Omega}u_{\varepsilon}v_{\varepsilon}
\end{align}
for all $t\in (0,T_{max,\varepsilon})$.

From \eqref{3.31}, \eqref{3.32} and \eqref{3.37}, there exists $C(p, K)>0$ such that
\begin{align}\label{3.38}
   &\frac{d}{dt}\left\{ \frac{1}{p}\int_{\Omega}u_{\varepsilon}^{p}+\int_{\Omega}\frac{|\nabla
   v_{\varepsilon}|^q}{v_{\varepsilon}^{q-1}}\right\}
   %+    \Gamma(q)\int_{\Omega}\frac{|\nabla    v_{\varepsilon}|^{q-2}}{v_{\varepsilon}^{q-3}}|D^2\ln{v_{\varepsilon}}|^2
   +\frac{\Gamma(q)}{4}\int_{\Omega}u_{\varepsilon}\frac{|\nabla
   v_{\varepsilon}|^q}{v_{\varepsilon}^{q-1}}+\frac{p-1}{4}\int_{\Omega}u_{\varepsilon}^{p-1}v_{\varepsilon}|\nabla
   u_{\varepsilon}|^2\nonumber\\
   &\leq C(p,K)\int_{\Omega}\left( v_{\varepsilon}|\nabla v_{\varepsilon}|^2+u_{\varepsilon}v_{\varepsilon}+v_{\varepsilon} \right).
\end{align}
 Integrating \eqref{3.38} with respect to time, we  have
\begin{align*}
    &\int_{\Omega}u_{\varepsilon}^{p}+\int_{\Omega}\frac{|\nabla
    v_{\varepsilon}|^q}{v_{\varepsilon}^{q-1}}+\frac{\Gamma(q)}{4}\int_0^{t}\int_{\Omega}u_{\varepsilon}\frac{|\nabla
    v_{\varepsilon}|^q}{v_{\varepsilon}^{q-1}}+\frac{p-1}{4}\int_0^{t}\int_{\Omega}u_{\varepsilon}^{p-1}v_{\varepsilon}|\nabla
    u_{\varepsilon}|^2
    \\
    &\leq C(p,K)\int_0^{t}\int_{\Omega}\left(v_{\varepsilon}|\nabla v_{\varepsilon}|^2+u_{\varepsilon}v_{\varepsilon}+v_{\varepsilon}\right)
      +\int_{\Omega}(u_{0}+1)^p+\int_{\Omega}\frac{|\nabla v_{0}|^q}{v_0^{q-1}}
\end{align*}
for all $t\in (0,T_{max,\varepsilon})$. Hence  from \eqref{2.11}, \eqref{2.18}, \eqref{1.6} and \eqref{2.10}, it
follows that
\begin{align}\label{3.39}
    \int_{\Omega}u_{\varepsilon}^{p}+\frac{p-1}{2}\int_0^{t}\int_{\Omega}u_{\varepsilon}^{p-1}v_{\varepsilon}|\nabla u_{\varepsilon}|^2\leq C(p,
    K).
\end{align}
Furthermore, \eqref{3.30-1} results from \eqref{3.37}, \eqref{3.39} and \eqref{2.10}, and hence completes the proof.
\end{proof}

\end{document}
